\section{Lecture 21 (November 19th)}
\begin{defi}
(Aharonov-Bohm effect) Consider a vector potential $\vec{A}$ and the generated magnetic field $\vec{B}=\nabla \times \vec{A}$. In the quantum world, the vector potential is more fundamental in some sense than the emergent field. Consider a solinoid of radius $a$ with a homogeneous magnetic field with strength $B_0$. What is the vector potential? Well, the answer is, from the homework,
\[\vec{A}=\begin{cases}
A_{\phi }(\rho <a)\sim \rho B_{0}/z\\
A_{\phi }(\rho >a)\sim \mathrm{\Phi} _{B}/2\pi \rho 
\end{cases}\]
with $\mathrm{\Phi} _{B}=B_0\pi a^2$. Let's try solving Schrodinger's equation outside the solenoid with $\vec{B}=0$ by trying to set $\vec{A}'=0$ outside. 
\[\vec{A}'=0=\vec{A}-\Big(\hat{\rho }\dfrac{\partial }{\partial \phi }+\hat{\phi }\dfrac{1}{\rho }\dfrac{\partial }{\partial \phi }+\hat{z}\dfrac{\partial }{\partial z}\Big)\mathrm{\Lambda } \]
we find that we need $\mathrm{\Lambda }=\mathrm{\Phi} \phi /2\pi $. We recall that
\[\psi' =\exp \Big[\dfrac{-iq\mathrm{\Lambda }}{\hbar }\Big]\psi =\exp \Big[\dfrac{ie}{\hbar }\dfrac{\mathrm{\Phi} }{2\pi }\phi \Big]\psi =\exp \Big[\dfrac{ie}{\hbar }\int \vec{A}\cdot d\vec{r}\Big]\psi \]
as $d\vec{r}=\rho \,d\vec{\phi }$. We of course need our single-valuedness condition,
\[\psi (\phi +2\pi )=\psi(\phi ) \]
but
\[\psi (\phi +2\pi )=e^{ie\mathrm{\Phi} /\hbar }\psi (\phi )\]
and the phase of the wavefunction changes. For a two slit interference experiment,
\[\begin{align*}
\psi_1+\psi_2=&S_1e^{iP_1}+S_2e^{iP_2}\\
=&e^{iP_1}(S_1+S_2e^{iP_2-iP_1})
\end{align*}\]
where the phase difference is affected by the nontrivial integral
\[\oint \vec{A}\cdot d\vec{r} \]
and the interference pattern changes. A related problem will be in the test! 
\end{defi}
\vspace{2ex}
\begin{thm}
(Radiation theory) First of all, define
\[H=\dfrac{(\vec{p}-(-e)\vec{A})^2}{2m_{e}}\]
where $\vec{A}$ is given as the vector potential of light. The most widely used is the coulomb gauge, with $\nabla \cdot \vec{A}=0$. In vaccum, $\mathrm{\Phi} =0$, and we simply have
\[\vec{A}=\nabla \times \vec{A}\qquad\&\qquad \vec{E}=-\dfrac{\partial \vec{A}}{\partial t} \]
Second of all, we must discuss whether our system is semi-classical or quantum. For a true quantum picture, we need 
\[a^{r}(k)\qquad\&\qquad a^{r}(k)^{\dagger }\qquad [a^{r}(k),a^{r}(k)^{\dagger }]=1\]
The field becomes, approximately
\[\vec{A}(x,t)=\sum _{k}\sum _{r=1}^{2}\Big[\vec{\epsilon }\,^{r}(k)a^{r}(k)e^{-i\omega t+ikx}+\vec{\epsilon }\,^{r}(k)^{\dagger }a^{r}(k)^{\dagger }e^{i\omega t-ikx}\Big]\]
Most simply, let's consider monochromatic light with just a single $\omega $. Noting that
\[N\hbar \omega =\int_{V} d^3\vec{x}\,\Big(\dfrac{\epsilon_0}{2}\vec{E}\cdot \vec{E}+\dfrac{\mu_0}{2}\vec{B}\cdot \vec{B}\Big)\]
and $\vec{E}$ and $\vec{B}$ from above, we can determine the normalisation of the ansatz above.
\[2\int d^3\vec{x}\,\dfrac{\epsilon_0}{2}\langle \vec{E}\cdot \vec{E}\rangle _{\mathrm{avg}}=2\omega ^2V\epsilon_0C^2\]
which gives
\[C=\sqrt{\dfrac{N\hbar  }{2\epsilon_0\omega V}}\]
\end{thm}
\vspace{2ex}

